\begin{table}[t]
\centering
\caption{Image-level bootstrap differences, reliability-aware $p=0.15$ minus matched early fusion. Values are descriptive percentile intervals from 2,000 image resamples.}
\label{tab:bootstrap}
\scriptsize
\begin{tabular}{@{}lcc@{}}
\toprule
Metric & Mean difference & 95\% CI\\
\midrule
AP50 & +0.0177 & [+0.0153, +0.0203]\\
AP75 & +0.0553 & [+0.0483, +0.0622]\\
F1 & +0.0187 & [+0.0150, +0.0226]\\
\bottomrule
\end{tabular}
\end{table}

\subsection{Synthetic Channel Removal}
To test controlled missing-input behavior, one channel group is deterministically set to zero at a time while weights, thresholds, and postprocessing remain fixed. \tabref{tab:removal} summarizes the archived seed-0/seed-2 checkpoint means. The reliability-aware system retains higher AP50 after removal of RGB, thermal, or event inputs. The largest separation follows thermal removal, where AP50 is 0.7030 for reliability-aware fusion and 0.3648 for early fusion. This deterministic intervention does not emulate measured physical sensor outages, temporal misalignment, or cross-device deployment failures.

\begin{table*}[t]
\centering
\caption{Synthetic channel removal on component-disjoint validation. Each value is the two-checkpoint mean. Delta is relative to the all-modal mean of the same model group. The intervention is deterministic zero-channel input, not a physical sensor-failure experiment.}
\label{tab:removal}
\scriptsize
\begin{tabular}{@{}llrrrrr@{}}
\toprule
Model group & Condition & Precision & Recall & F1 & AP50 & AP75\\
\midrule
Matched early & All modalities & 0.9188 & 0.8720 & 0.8945 & 0.9408 & 0.8208\\
Matched early & RGB removed & 0.8202 & 0.6531 & 0.7269 & 0.7379 & 0.5265\\
Matched early & Thermal removed & 0.8226 & 0.3003 & 0.4400 & 0.3648 & 0.2626\\
Matched early & Event removed & 0.9001 & 0.8731 & 0.8863 & 0.9339 & 0.8021\\
\midrule
Reliability $p=0.15$ & All modalities & 0.9273 & 0.8997 & 0.9133 & 0.9586 & 0.8760\\
Reliability $p=0.15$ & RGB removed & 0.8569 & 0.8229 & 0.8394 & 0.8985 & 0.7051\\
Reliability $p=0.15$ & Thermal removed & 0.8321 & 0.5885 & 0.6894 & 0.7030 & 0.3883\\
Reliability $p=0.15$ & Event removed & 0.9305 & 0.8970 & 0.9134 & 0.9582 & 0.8600\\
\bottomrule
\end{tabular}
\end{table*}

\begin{figure}[t]
\centering
\includegraphics[width=\columnwidth]{figures/fig4_channel_removal.pdf}
\caption{AP50 under deterministic synthetic channel removal. Reliability-aware $p=0.15$ retains higher AP50 under all four input conditions.}
\label{fig:removal}
\end{figure}

\subsection{Checkpoint-Backed Single-Modality Replication}
The original checkpoints underlying the earlier author-supplied single-modality table were unavailable. We therefore conducted a fresh V81 replication rather than treating those unidentified records as recoverable evidence. RGB-only, thermal-only, and event-only detectors were retrained for seeds 0, 1, and 2 under the frozen component-disjoint protocol: 50 epochs, batch size 4, $640\times640$ inputs, AdamW at $10^{-4}$ with weight decay $10^{-4}$, no modality dropout, and selection of the checkpoint with the highest project-local development-validation AP50. No guard access, tuning, early stopping, seed replacement, selective rerun, or checkpoint substitution was used.

All nine retained checkpoints were then evaluated once with the standardized COCO contract used for the primary TriAir analysis: score threshold 0.001, NMS threshold 0.6, and at most 100 detections per image. Every result is linked to its selected epoch, checkpoint SHA256, and a common validation-manifest identity (SHA256 prefix \texttt{722efc6f74a7615a}). Full hashes and evaluator JSON records are archived with the project. Because V81 is a fresh replication, its values are not presented as recovered identities for the earlier supplied table, which is retained only as a provenance reconciliation record.

\begin{table*}[t]
\centering
\caption{Checkpoint-backed V81 single-modality replication on component-disjoint TriAir development-validation. AP denotes COCO AP@[0.50:0.95]. Epoch is the retained checkpoint selected by project-local AP50.}
\label{tab:single-modal-v81-seeds}
\scriptsize
\setlength{\tabcolsep}{2.8pt}
\resizebox{\textwidth}{!}{%
\begin{tabular}{@{}lrrrrrrrr@{}}
\toprule
Modality & Seed & Epoch & AP & AP50 & AP75 & AR1 & AR10 & AR100\\
\midrule
RGB-only & 0 & 12 & 0.4508 & 0.7638 & 0.4521 & 0.1642 & 0.5258 & 0.5892\\
RGB-only & 1 & 7 & 0.4467 & 0.7676 & 0.4438 & 0.1648 & 0.5231 & 0.5922\\
RGB-only & 2 & 9 & 0.4443 & 0.7709 & 0.4325 & 0.1661 & 0.5187 & 0.5876\\
Thermal-only & 0 & 5 & 0.5256 & 0.8156 & 0.5864 & 0.2077 & 0.5868 & 0.6552\\
Thermal-only & 1 & 2 & 0.4977 & 0.8340 & 0.5500 & 0.1941 & 0.5661 & 0.6321\\
Thermal-only & 2 & 2 & 0.5356 & 0.8463 & 0.5965 & 0.2087 & 0.5948 & 0.6546\\
Event-only & 0 & 4 & 0.1942 & 0.3627 & 0.1956 & 0.0785 & 0.2704 & 0.3571\\
Event-only & 1 & 3 & 0.1962 & 0.3654 & 0.1889 & 0.0747 & 0.2701 & 0.3617\\
Event-only & 2 & 4 & 0.1941 & 0.3690 & 0.1984 & 0.0719 & 0.2679 & 0.3486\\
\bottomrule
\end{tabular}%
}
\end{table*}

\begin{table*}[t]
\centering
\caption{V81 standardized single-modality COCO means and sample standard deviations over three fresh training seeds.}
\label{tab:single-modal-v81-summary}
\scriptsize
\setlength{\tabcolsep}{2.7pt}
\resizebox{\textwidth}{!}{%
\begin{tabular}{@{}lrrrrrr@{}}
\toprule
Modality & AP & AP50 & AP75 & AR1 & AR10 & AR100\\
\midrule
RGB-only & $0.4473\pm0.0033$ & $0.7674\pm0.0036$ & $0.4428\pm0.0098$ & $0.1650\pm0.0009$ & $0.5225\pm0.0036$ & $0.5897\pm0.0024$\\
Thermal-only & $\mathbf{0.5196\pm0.0196}$ & $\mathbf{0.8320\pm0.0154}$ & $\mathbf{0.5776\pm0.0244}$ & $\mathbf{0.2035\pm0.0081}$ & $\mathbf{0.5826\pm0.0148}$ & $\mathbf{0.6473\pm0.0132}$\\
Event-only & $0.1949\pm0.0012$ & $0.3657\pm0.0032$ & $0.1943\pm0.0049$ & $0.0751\pm0.0033$ & $0.2694\pm0.0014$ & $0.3558\pm0.0067$\\
\bottomrule
\end{tabular}%
}
\end{table*}

Thermal is the strongest standalone modality for every standardized mean, while event-only is the weakest. Because the all-modal V48 checkpoints and the V81 replication use the same component-disjoint validation manifest, standardized COCO metric definitions, and seed labels, compatible seed-paired contrasts can be computed without mixing evaluator contracts (\tabref{tab:fusion-over-v81-thermal}). The full reliability-aware system exceeds thermal-only by $0.1960\pm0.0329$ AP, $0.1215\pm0.0154$ AP50, and $0.2952\pm0.0349$ AP75; each contrast is positive for all three seeds. Matched early fusion also exceeds thermal-only for every seed. These descriptive comparisons support a multimodal benefit beyond the strongest individual stream under the frozen development protocol, but do not establish statistical significance or physical sensor-failure robustness.

\begin{table*}[t]
\centering
\caption{Seed-paired standardized COCO contrasts against the authoritative V81 thermal-only replication. Positive pairs count the three shared seed labels.}
\label{tab:fusion-over-v81-thermal}
\scriptsize
\begin{tabular}{@{}llrrr@{}}
\toprule
Multimodal system minus thermal-only & Metric & Mean delta & Sample SD & Positive pairs\\
\midrule
Reliability-aware $p=0.15$ & AP & +0.1960 & 0.0329 & 3/3\\
Reliability-aware $p=0.15$ & AP50 & +0.1215 & 0.0154 & 3/3\\
Reliability-aware $p=0.15$ & AP75 & +0.2952 & 0.0349 & 3/3\\
Matched early fusion & AP & +0.1606 & 0.0223 & 3/3\\
Matched early fusion & AP50 & +0.1053 & 0.0199 & 3/3\\
Matched early fusion & AP75 & +0.2314 & 0.0265 & 3/3\\
\bottomrule
\end{tabular}
\end{table*}

